\section{Friction}
In high school, we deal with two types of friction: static and kinetic.

\subsection{Kinetic Friction}
Use this when the two surfaces in contact \textit{slide relative to each other}.
\begin{equation}
    F_{fk} = \mu_{fk} * F_{n}
\end{equation}

Kinetic friction on an object can cause the object to \underline{slow down} or \underline{speed up}.

\subsection{Static Friction}
Use this when the two surfaces are \underline{not sliding} relative to each other.
\begin{equation}
    F_{fs.max} = \mu_{fs} * F_{n}
\end{equation}

\subsection{Remainder}
Do not forget to justify \underline{$F_{N} = F_{g}$}.
\newpage
\mypic{pictures/example9.png}{}{1}
